\documentstyle[% 


righttag% 


,11pt% 


,amssymb% 


,russian% 


,izvru% 


]{amsart} 


 


%       Math definitions 


 


\newcommand{\A}{{\cal A}} 


%\newcommand{\div}{\operatorname{div}} 


\newcommand{\diag}{\operatorname{diag}} 





%\theoremstyle{plain} %% This is the default


%\begingroup


%\theorembodyfont{\it}


%\newtheorem{lemnonum}{{\hskip\parindent Лемма}}


%\renewcommand{\thelemnonum}{}


%\endgroup


 


\setcounter{page}{77} 


\god{1997} 


\nomer{4~(419)} 


\udk{519.63} 


\author{\it Л.А.\,КРУКИЕР} 


\title{КОСОСИММЕТРИЧНЫЕ ИТЕРАЦИОННЫЕ МЕТОДЫ РЕШЕНИЯ\\ 


СТАЦИОНАРНОЙ ЗАДАЧИ КОНВЕКЦИИ-ДИФФУЗИИ\\ 


С МАЛЫМ ПАРАМЕТРОМ ПРИ СТАРШЕЙ ПРОИЗВОДНОЙ} 


 


%\thanks{Работа выполнена при поддержке РФФИ, грант \No. 94-01-00183.} 


\date{} 


\pagestyle{plain} 


\begin{document} 


 


\maketitle 


 


\section{Введение} 


 


Наличие малого параметра при старшей производной в уравнении 


конвекции-диффузии эквивалентно преобладанию конвективных членов уравнения 


над диффузионными и сильно осложняет численное решение. При определенных 


условиях такая ситуация соответствует возникновению пограничного слоя, т.е. 


сильному росту решения на узком участке области расчета. Это явление 


характерно для широкого круга решаемых задач \cite{Lan}. Аппроксимация 


подобных задач конечными разностями приводит к необходимости решать сильно 


несимметричные системы линейных алгебраических уравнений. 


 


Под сильно несимметричными системами линейных алгебраических уравнений 


\begin{equation} 


Au=f, \quad A=A_0+A_1,\quad A_0=A_0^*,\quad A_1=-A_1^*, 


\end{equation} 


будем понимать системы, у которых симметричная часть матрицы гораздо меньше 


(в смысле некоторой нормы), чем ее кососимметричная часть. 


 


Однако наличие малого параметра при старшей производной является 


необходимым, но не достаточным условием возникновения пограничного слоя в 


решаемой задаче. Появление пограничного слоя существенно зависит от краевых 


условий и правой части исходных дифференциальных уравнений, т.е. от свойств 


функции $f$ в (1), которая после аппроксимации дифференциального уравнения 


содержит эти данные. 


 


Рассмотрим в области $\Omega$ следующую краевую задачу: 


\begin{equation} 


-\frac{1}{Pe}\Delta U+v_1\frac{\partial U}{\partial x}+ 


v_2\frac{\partial U}{\partial y}=f(x,y),\qquad 


\left. U\right|_{\partial \Omega}=U_{\text{гр}}. 


\end{equation} 


Наиболее сложный этап аппроксимации задачи (2) связан с выбором аппроксимации 


для конвективных членов (первых производных) уравнения, хотя число вариантов 


достаточно ограничено, т.к. используются либо аппроксимации первого порядка 


(разности ``вперед'' или ``назад''), либо центрально-разностная аппроксимация 


2 порядка точности. 


 


Известно \cite{Ro}, что при использовании разностей ``вперед'' или ``назад'' 


необходимо сохранять в разностной форме свойства монотонности и принцип 


максимума, которые при преобразовании разностной схемы в систему линейных 


алгебраических уравнений дают в (1) $M$-матрицу \cite{Neu}. (Если знаки 


функций $v_1(x,y)$ и 


$v_2(x,y)$ непостоянны, то вместо разностей ``вперед'' или ``назад'' 


используются разности 


``против потока''.) Для этого случая можно предложить 


достаточно широкий набор 


итерационных методов, эффективно решающих системы с $M$-матрицами 


\cite{Y}, и, 


в первую очередь, методы, основанные на неполном разложении Холецкого в 


сочетании с градиентными методами \cite{Mei}. 


Но низкая точность схем первого 


порядка (особенно в случае сильно меняющихся коэффициентов задачи) и ошибки 


при описании поведения решения в пограничном слое не позволяют рекомендовать 


их для использования во всех случаях. 


 


Центрально-разностные схемы, дающие более высокий порядок аппроксимации 


и правильно описывающие поведение решения в пограничном слое, накладывают 


достаточно жесткие ограничения на шаг сетки. Это связано с необходимостью 


сохранения диагонального преобладания в исходной матрице системы (1), которое 


достаточно для сходимости большинства базовых и наиболее распространенных 


итерационных методов таких, как методы Якоби, Зейделя, SOR и SSOR. В случае 


отсутствия диагонального преобладания в исходной матрице часть этих методов 


(напр., метод Зейделя) перестает сходиться, а другие (напр., SOR) сильно 


уменьшают скорость сходимости. 


 


Очевидно, наличие в уравнении малого параметра при старшей производной 


накладывает существенные ограничения на шаги по пространству 


центрально-разностной схемы, если мы хотим сохранить в получаемой матрице 


$A$ 


системы (1) диагональное преобладание. Такие ограничения увеличивают объем 


вычислительной работы, необходимую память и время вычислений. Вместе с тем, 


они не связаны с существом решаемой задачи, а накладываются итерационным 


методом, т.к. в большей части области расчета решение сильно не меняется, что 


позволяет считать с достаточно большим шагом по пространству, сохраняя 


приемлемую точность. 


 


Предлагаемый ниже класс итерационных методов позволяет избавиться от 


таких ограничений, т.к. его использование требует от исходной несимметричной 


матрицы лишь положительной определенности ее симметричной части, а не 


диагонального преобладания в исходной матрице. Таким образом, предлагаемый 


класс итерационных методов позволяет эффективно решать сильно несимметричные 


системы, полученные из (2) при числах Пекле $Pe\ge 10^5.$ 


 


\section{Постановка задачи} 


 


Стационарная задача конвекции-диффузии в несжимаемой среде описывается 


уравнением (2) с добавлением уравнения неразрывности 


\begin{equation} 


{\mathrm{div}}\,V=0,\quad V=\{v_1,v_2\}. 


\end{equation} 


Задача (2) записана в недивергентной форме, но учитывая уравнение (3) 


несжимаемости среды, ее можно переписать в виде 


\begin{equation} 


-\frac{1}{Pe}\Delta U + 


\frac{\partial v_1U}{\partial x}+


v_2\frac{\partial v_2U}{\partial y}=f(x,y),\qquad 


\left. U\right|_{\partial \Omega}=U_{\text{гр}} 


\end{equation} 


или в так называемом ``симметричном'' виде \cite{Vab}, \cite{Kru1} 


\begin{equation} 


-\frac{1}{Pe}\Delta U + 


\frac{\partial v_1U}{\partial x}+


v_1\frac{\partial U}{\partial x}+ 


\frac{\partial v_2U}{\partial y}+


v_2\frac{\partial U}{\partial y} 


=f(x,y),\qquad 


\left. U\right|_{\partial \Omega}=U_{\text{гр}}.


\end{equation} 


При аппроксимации задач (2), (4) или (5) в области $\Omega$ с границей 


$\partial\Omega$


конечными разностями или конечными элементами получается система линейных 


алгебраических уравнений (СЛАУ) с несамосопряженной матрицей, причем свойства 


этой матрицы существенно зависят от способа аппроксимации первых производных 


и вида исходной задачи. Так, применяя к задаче (2) разности ``против 


потока'', получим в результате $M$-матрицу \cite{Ro}, а используя центральные 


разности в уравнении (5), --- положительную матрицу \cite{Kru1}. При этом 


учитывается уравнение несжимаемости среды (3), благодаря которому можно 


получить запись задачи в ``симметричном'' виде. 


 


Всюду в дальнейшем будем рассматривать в качестве исходной краевую 


задачу (5), для аппроксимации которой применяется конечно-разностная схема с 


центральными разностями. 


 


\section{Конечно-разностная аппроксимация уравнения конвекции-диффузии 


\newline 


и свойства получаемой СЛАУ} 


 


Введем в области $\Omega$ прямоугольную равномерную сетку $M$ с шагами 


$h_1$ и $h_2$ соответственно по направлению $Ox$ и $Oy$. Краевые условия на 


$\partial\Omega$ интерполируются со 


вторым порядком точности на границу $\Gamma$ сеточной области $M$. 


Запишем краевую 


задачу (5) на сетке $M$, используя стандартные обозначения из \cite{Sam} 


для ячейки с индексами $(i,j)$. Тогда 


\begin{equation} 


-\frac{1}{Pe}\Delta_h u+v_1 u_{\overset\circ x} 


+(v_1 u)_{\overset\circ x}+v_2 u_{\overset\circ y}+ 


(v_2 u)_{\overset\circ y}=f_{ij},\qquad \left. u_{ij}\right|_\Gamma 


=u_{\Gamma pij}. 


\end{equation} 


Запишем разностную задачу (6) на стандартном пятиточечном шаблоне, учитывая 


краевые условия в правой части получаемой системы и умножая все уравнение на 


$h^2$ и $Pe$ (в предположении, что $h_1=h_2=h$) 


\begin{gather} 


-u_{i-1\,j}-u_{i+1\,j}+4u_{ij}-u_{i\,j-1}-u_{i\,j+1}+


k[(v_{1\,ij}+v_{1\,i+1j})u_{i+1\,j}-(v_{1\,i-1j}+v_{1\,ij}) 


u_{i-1\,j}+\notag\\


+(v_{2\,ij+1}+v_{2\,ij})u_{i\,j+1}- 


(v_{2\,ij}+v_{2\,ij-1})u_{i\,j-1}]=F_{ij}. 


\end{gather} 


 


Здесь величина 


\begin{equation} 


k=Pe\cdot h/2 


\end{equation} 


называется коэффициентом кососимметрии задачи. 


 


В результате получим систему линейных несамосопряженных уравнений (1). 


 


Известно \cite{Sam}, что любой линейный оператор (матрицу) можно 


представить в виде суммы симметричного и кососимметричного операторов 


(матриц). В частности, $A=A_0+A_1,$ где $A_0$ --- симметричная, а $A_1$ --- 


кососимметричная части матрицы $A$. 


 


\begin{defn} Оператор $A$ называется положительным, если его симметричная 


часть положительно определена. 


\end{defn} 


\begin{thm} Оператор $A$ системы {\em (1),} построенный для уравнения 


{\em (5)} по 


разностной схеме {\em (7),} несамосопряжен и положителен. 


\end{thm} 


Доказательство этой теоремы следует из того, что разностный оператор (7) 


задачи (5) явным образом представим в виде суммы симметричного положительно 


определенного оператора $A_0$, 


являющегося разностным аналогом оператора Лапласа, 


и кососимметричного оператора $A_1$, 


являющегося разностным аналогом конвективных 


членов уравнения (5). 


 


Отметим, что запись уравнения конвекции-диффузии в виде (5) позволяет 


при центрально-разностной аппроксимации первых производных получить сразу 


кососимметричный оператор для любых коэффициентов уравнения. Вместе с тем 


необходимо отметить, что лишь в случае постоянных коэффициентов уравнение 


конвенции-диффузии, записанное в виде (2) или (4), при центрально-разностной 


аппроксимации дает сразу кососимметричный оператор.  


 


Особую роль при 


решении системы (1) играет коэффициент кососимметрии (8), т.к. его значение 


определяет возможность использования различных итерационных методов. Так при 


$k\le 1$ матрица системы (1) обладает свойством диагонального преобладания 


(предполагается, что скорости течения нормированы и не превосходят по модулю 


1) и, следовательно, для решения системы можно использовать хорошо известные 


и эффективные методы \cite{Sam}--\cite{Ha}, 


упомянутые ранее. При $k\gg 1$ эти методы работают 


плохо или вообще расходятся, и необходим другой метод, использующий 


особенности задачи. 


 


\section{Итерационный метод} 


 


Для решения СЛАУ (1) с матрицей $A$ используем стационарный двуслойный 


итерационный метод вида 


\begin{equation} 


B\frac{y_{j+1}-y_j}{\tau}+Ay_j=f,\quad j=1,2,\dotsc , 


\end{equation} 


где $B,$ $A$ --- линейные операторы в $H,$ $H$ --- конечномерное гильбертово 


пространство, $y_0, f\in H,$ $\tau$ --- параметр, $\tau>0,$ $B$ --- обратимый 


оператор, $y_j$ --- $j$-e приближение, $y_0$ --- начальное приближение. 


 


Заметим, что для матрицы $A$ справедливы следующие равенства: 


$$ 


A=A_0+A_1;\quad A_0=1/2(A+A^*)=A^*_0;\quad A_1=1/2(A-A^*)=-A_1^*;\quad 


A_1=K_B+K_H;\quad K_B=-K^*_H, 


$$ 


где $A_0$ --- симметричная, $A_1$ --- кососимметричная части $A,$ 


$K_H$ и $K_B$ --- строго нижняя или строго 


верхняя треугольные части матрицы $A_1$ соответственно. 


 


Не нарушая общности, считаем все диагональные элементы матрицы $A$ равными 


1, т.е. система (1) масштабирована. 


 


Оператором перехода метода (9) является 


\begin{gather*} 


G=B^{-1}(B-\tau A)=(B_0+B_1)^{-1}(B_0+B_1 -\tau A_0 -\tau A_1);\\


B_0=1/2(B+B^*)=B^*_0,\quad B_1=1/2(B-B^*)=-B_1^*. 


\end{gather*} 


Предположим, что 


\begin{equation} 


B_0=B_0^* >0. 


\end{equation} 


Тогда 


\begin{gather*} 


G=[B^{1/2}_0(E+B^{-1/2}_0B_1\overline B^{1/2}_0)B_0^{1/2}]^{-1} 


[B^{1/2}_0(E+B^{-1/2}_0B_1B^{-1/2}_0-\tau B^{-1/2}_0A_0B^{-1/2}_0 


-\tau B^{-1/2}_0 A_1B^{-1/2}_0)]=\\ 


=B^{1/2}_0(E+B^{-1/2}_0 B_1 B^{-1/2}_0)^{-1} 


(E+B^{-1/2}_0 B_1 B^{-1/2}_0-\tau B^{-1/2}_0 A_0 B^{-1/2}_0 


-\tau B^{-1/2}_0 A_1 B^{-1/2}_0)B^{1/2}_0. 


\end{gather*} 


Введем операторы 


\begin{equation} 


P_0=B^{1/2}_0 A_0 B^{-1/2}_0 = P_0^* >0,\quad 


P_1 = B^{-1/2}_0 B_1 B^{-1/2}_0 = -P_1^*,\quad 


P_2 = B^{-1/2}_0 A_1 B^{-1/2}_0=-P_2^*. 


\end{equation} 


Запишем 


\begin{equation} 


G=B^{-1/2}_0 L B^{1/2}_0, 


\end{equation} 


где 


\begin{equation} 


L=(E+P_1)^{-1}(E+P_1-\tau P_0 -\tau P_2). 


\end{equation} 


Будем изучать сходимость метода (9) в энергетической норме 


$\|x\|_{B_0}=(x,x)^{1/2}_{B_0}=(B_0 x, x)^{1/2}$. Используем в 


(9) замену переменных $y=B^{-1/2}_0 x$. Тогда из (12) получим 


\begin{equation} 


\|G\|_{B_0}=\|L\|. 


\end{equation} 


Поскольку вид оператора $B$ еще не оговаривался, кроме условия (10) его 


положительной определенности, потребуем, чтобы оператор  


$(E+P_1 -\tau P_0 - \tau P_2)$ в (13) был 


самосопряжен. Этого можно добиться с помощью соотношения 


$$ 


P_1 = \tau P_2, 


$$ 


из которого после подстановки выражения (11) для $P_1$ и $P_2$ получим 


$$ 


B^{-1/2}_0 B_1 B^{-1/2}_0=\tau B^{-1/2}_0 A_1 B^{-1/2}_0, 


$$ 


что эквивалентно равенству 


\begin{equation} 


B_1=\tau A_1. 


\end{equation} 


Это равенство является основным при дальнейшем конструировании итерационных 


методов для сильно несимметричных задач. 


 


С учетом (12)--(15) получается 


 


\begin{thm} Для сходимости итерационного метода {\em (9)} с оператором  


$B,$ удовлетворяющим условиям {\em (10)} и {\em (15)} в энергетическом 


пространстве $H_{B_0},$ достаточно, чтобы 


\begin{equation} 


\|(E+\tau P_1)^{-1}(E-\tau P_0)\| <1 


\end{equation} 


или 


\begin{equation} 


B_0 \ge \tau /2\,A_0>0. 


\end{equation} 


\end{thm} 


 


Необходимо отметить, что условие (16) неконструктивно и трудно 


проверяемо. Поэтому упростим его. Применяя неравенство треугольника для (16), 


получим 


$$ 


\|G\|_{B_0}=\|L\|=\|(E+\tau P_1)^{-1}(E-\tau P_0)\| \le 


\|(E+\tau P_1)^{-1}\|\cdot \|E-\tau P_0\| \le \|E-\tau P_0\| 


$$ 


с использованием неравенства $\|(E+\tau P_1)^{-1}\|\le 1$ 


в силу кососимметрии оператора 


$P_1$. Условие (16) будет обеспечено, если 


\begin{equation} 


\|(E-\tau P_0)\| < 1. 


\end{equation} 


Учтя самосопряженность оператора $P_0$ 


и свойства операторных неравенств \cite{Gul}, из 


(18) имеем 


$$ 


-E < E-\tau P_0 <E, 


$$ 


что эквивалентно неравенствам 


$$ 


2E > \tau P_0 >0. 


$$ 


Подставляя в них выражение для $P_0$ из (11), получаем (17). 


 


Рассмотрим возможные способы выбора оператора $B$, чтобы обеспечить  


 


а) сходимость итерационного метода (9); б) эффективность обращения 


оператора $B$. 


 


Чтобы выполнить требование а), достаточно, как это было показано выше, 


выполнения операторного неравенства (17). В работе \cite{Kru2} 


было показано, что 


для выполнения условия (15) оператор $B$ 


должен иметь определенную структуру, а 


именно 


\begin{equation} 


B=B_c+\tau ((1+j)K_H +(1-j)K_B),\quad j=\pm 1,\quad B_c=B_c^*. 


\end{equation} 


 


Оператор $B_c$ пока произволен, но он должен обеспечивать выполнение (10). 


Отметим, что  


\begin{equation} 


B_0 =B_c+\tau j(K_B - K_H),\quad j=\pm 1;\quad B_1=\tau A_1. 


\end{equation} 


 


Рассмотрим оператор $B_0$ более подробно. Имеет место 


\bigskip





{\bf Лемма.} {\it Пусть матрица $Z\ne 0$ с элементами $\{z_{ij}\}^n_1$


имеет вид}


$$ 


Z=k+K^*, 


$$ 


{\it где $K$ --- вещественная строго верхняя или нижняя треугольная матрица.


Тогда спектр $\sigma(Z)$ матрицы $Z$ вещественен, состоит из положительных, 


отрицательных и нулевых собственных чисел, и}


$$ 


-v \le \sigma(Z) \le v,\qquad v=\max_i|R_i|, \qquad 


R_i=\sum^N_{j=1}|z_{ij}|. 


$$ 


 


{\it Более того, если матрица $Z$ двуциклическая \cite{Gul}, 


то ее собственные числа 


расположены симметрично относительно начала координат и}


$$ 


v=\rho(Z)=\|Z\|_2, 


$$ 


{\it где $\rho(Z)$ --- спектральный радиус матрицы $Z,$ а $\|\cdot\|_2$ --- 


спектральная норма.}


\pagebreak% !!!! remove in eng.





Доказательство первой части этой леммы дано в \cite{Kru3}. 


Вторая часть следует 


из результатов, приведенных в \cite{Ha}. 


 


Отметим сразу же, что рассматриваемая матрица системы (1), полученная в 


результате стандартной пятиточечной аппроксимации уравнения 


конвекции-диффузии, является согласно результатам из \cite{Ha} 


двуциклической. 


 


Используем эту лемму для изучения спектра симметричной части $A_0$ матрицы 


$A.$ Эта матрица масштабирована, т.е. по главной диагонали у нее стоят единицы 


и она также, как основная матрица, является двуциклической. 


 


\begin{thm} Пусть матрица $A$ симметрична, положительно определена, 


масштабирована и является двуциклической. Тогда ее спектр 


\begin{equation} 


\sigma(A) \in (0,2). 


\end{equation} 


\end{thm} 


 


{\bf Доказательство.} Пусть $\lambda_k \in \sigma(A)$. 


Тогда для всех $\lambda_k$ из спектра 


матрицы $A$ имеем 


\begin{equation} 


\lambda_k (A)=\lambda_k(E+A_R) = 1+\lambda_k(A_R)>0, 


\end{equation} 


где $A_R=U+U^*,$ $U$ --- строго 


нижняя (верхняя) треугольная часть матрицы $A_R$. Тогда из 


предыдущей леммы следует, что собственные числа матрицы $A_R$ симметричны 


относительно начала координат, максимальное и минимальное собственные числа 


связаны равенством $\lambda_{\text{min}}(A_R)=-\lambda_{\text{max}}(A_R)$. 


В силу этого равенства из (22) следует (21). 


 


Для нахождения оптимального итерационного параметра необходима 


информация об интервале, на котором его нужно искать. Левая граница этого 


интервала известна и равна 0, а правую найдем из достаточного условия 


сходимости. Учитывая (21) и то, что симметричная часть матрицы $A_0$ 


системы (1) 


удовлетворяет условиям теоремы 3, перепишем достаточное условие (17) в виде 


\begin{equation} 


B_0>\tau E >0. 


\end{equation} 


Подставив в (23) выражение для $B_0$ из (20), получим 


\begin{equation} 


B_c+\tau jK>\tau E>0,\quad K=K_B-K_H. 


\end{equation} 


Так как оператор $jK$ имеет в силу леммы как положительные, так и 


отрицательные собственные числа и, кроме того, $\tau >0,$ 


то для выполнения (24) 


оператор $B_c$ должен быть положительным. 


 


Предположив, что оператор $B_c$ не зависит от параметра $\tau$ 


и рассмотрев лишь левое неравенство в (24) (правое неравенство выполняется 


всегда при $\tau >0),$ получим 


$$ 


B_c>\tau (E-jK)>0. 


$$ 


 


Введя $\lambda_{\text{min}}=\min\limits_i \lambda_i(B_c)$ 


и учитывая лемму для 


оператора $jK,$ получаем 


$$ 


\lambda_{\text{min}}>\tau (1+\rho(K))=\tau(1+\|K\|_2) \Longrightarrow 


\tau < \lambda_{\text{min}}/(1+\|K\|_2). 


$$ 


 


Отметим также два обстоятельства: $C$-нормы матриц $K$ и $A_1$ 


совпадают, т.к. 


эти матрицы имеют одинаковые по модулю элементы, и в силу нормальности этих 


матриц имеем из \cite{Gul} 


$$ 


\|K\|_2 \le \|K\|_\infty =\|A_1\|_\infty. 


$$ 


 


При выборе параметра $\tau$ удобней, с точки зрения эффективности вычислений, 


использовать именно $C$-норму. Тогда верхняя граница интервала для 


итерационного параметра, при котором метод (9) будет сходиться, определяется 


формулой 


\begin{equation} 


\tau _*=\frac{\lambda_{\text{min}}} 


{1+\|A_1\|_\infty}. 


\end{equation} 


 


Таким образом, доказано следующее достаточное условие сходимости метода. 


 


\begin{thm} Пусть матрица $A$ системы {\em (1)} положительна. Тогда 


итерационный метод {\em (9)} с оператором $B$, заданным формулой {\em (19),} и 


с оператором $B_c,$ не зависящим от итерационного параметра $\tau ,$ сходится 


в пространстве $H_{B_0}$ при $\tau\in(0,\tau_*),$ где $\tau_*$ определено в 


{\em (25).} 


\end{thm} 


Рассмотрим несколько способов задания оператора $B_c$. Очевидно, что 


наиболее эффективным из них в смысле вычислительной работы является оператор 


$B_c$ с диагональной структурой элементов. Это приводит к треугольной 


структуре оператора $B,$ что позволяет его легко обращать. 


 


В частности, в работе \cite{Kru3} 


\begin{equation} 


B_c=E, 


\end{equation} 


что обеспечивало выполнение условий предыдущих теорем. Там же были приведены 


результаты, связанные с оптимальным выбором итерационного параметра, и дана 


оценка скорости сходимости метода. 


 


Вместе с тем, проведенные числовые эксперименты показывают, что 


предложенные методы требуют большого числа итераций в случаях, когда скорости 


$v_1$ и $v_2$ сильно меняются. Это связано с тем, что операторы верхнего слоя 


из \cite{Kru3} не содержат достаточной информации об изменении коэффициентов 


задачи. 


 


Чтобы устранить этот недостаток, в \cite{Kru4} предложено использовать 


оператор 


\begin{equation} 


B_c = E+\omega D, 


\end{equation} 


где $D=\diag(K_H K_B+K_B K_H)/2.$ В этом случае оператор $B_c$ уже содержит 


информацию об изменениях 


коэффициентов в уравнении (5) и адекватно на это реагирует, что и 


подтверждается расчетами. Вместе с тем, теоретические оценки сходимости 


метода с $B_c$ из (27) совпадают с оценками для метода с $B_c$ из (26), что 


указывает на грубость используемых оценок. 


 


\section{Численные результаты} 


 


Было проведено численное исследование данного класса итерационных 


методов на примере нахождения решения системы линейных алгебраических 


уравнений, полученной в результате разностной аппроксимации (6) задачи (5) 


при следующих значениях коэффициентов: 


\medskip 





%{\hspace{13cm} Таблица 1} 


{\hfill Таблица 1\hfill} 


 


%{\hfill Значения коэффициентов уравнения (5)\hfill} 





\par


\bigskip


{\hfill


{\begin{tabular}{|c|c|c|} 


\hline 


Задача & $v_1$ & $v_2$ \\ 


\hline 


1 & 1 & -1 \\ 


\hline 


2 & $1-2x$ & $2y-1$ \\ 


\hline 


3 & $x+y$ & $x-y$ \\ 


\hline 


4 & $\sin 2\pi x$ & $2\pi y \cos 2\pi x$\\ 


\hline 


\end{tabular} 


}


\hfill}


\par 


\bigskip


В качестве точного решения уравнения (5) в области $[0;1]\times [0,1]$ 


берется степенная 


функция $u=x^m + y^m,$ где $m=5,$ 


моделирующая решение типа пограничного слоя, и функция \linebreak


% --------------------------------------- join in eng.  ^^^^^^^^^


$u=\sin \pi x \sin \pi y,$ 


моделирующая гладкое решение. 


 


Итерационный процесс прекращается, если 


$$ 


R=\|r^{(k)}\|_2 / \|r^{(0)}\|_2 \le \varepsilon, 


$$ 


где $r^{(k)},$ $r^{(0)}$ --- невязки на $k$-й и на 0-й итерации, 


$\varepsilon=10^{-6}.$ 


 


В приведенной таблице 2 содержатся значения числа итераций, необходимых 


для достижения указанной точности при решении задачи на сетке $32\times 32$ 


в случае 


использования оператора $B_c$, заданного формулой (26) --- ТМ(1)Г и ТМ(1)П и 


формулой (27) --- ТМ(D)Г и ТМ(D)П. Буквы Г и П после скобок означают, что в 


качестве точного решения задачи использовалось, соответственно, гладкое 


решение или решение типа пограничного слоя. 


 


Из приведенных данных видно, что тип решения не оказывает заметного 


влияния на число итераций, т.е. сходимость метода не зависит от поведения 


решения, а связана со значением малого параметра при старшей производной и 


поведением скоростей течения. Уменьшение шага по пространству сказывается


лишь на точности получаемого решения. Теоретически поведение методов 


(9), (26) и (9), (27) не отличается. Численно проверено, что скорость 


сходимости метода (9), (27) примерно в два раза выше лишь в задаче 4,где 


характеристики течения сильно менялись. В остальных случаях метод (26) 


оказался эффективней, так как требует на 25\% меньше арифметических действий за 


счет простой структуры оператора $B.$ 


\medskip 





%{\hspace{13cm} Таблица 2} 


{\hfill Таблица 2\hfill} 


 


{\hfill Результаты решения задач на сетке $32\times 32$ \hfill} 





\par


\bigskip


{\hfill


{\begin{tabular}{|c|c|c|c|c|c|} 


\hline 


Номер задачи & $Pe$ & ТМ(1)Г & ТМ(1)П & ТМ(D)Г & ТМ(D)П\\ 


\hline 


1  &  $10^3$  &  216 &      214 &      214  &     219 \\ 


\hline 


2 &   $10^3$ &   336 &     335 &      288 &     263 \\ 


\hline 


3 &   $10^3$ &   305 &     299 &     264 &     253 \\ 


\hline 


4 &   $10^3$  &  606 &     612 &      341 &     360 \\ 


\hline 


1 &   $10^4$  &  1517 &    1528 &    1486 &    1394 \\ 


\hline 


2 &   $10^4$ &   1462 &    1403 &    1153 &   994 \\ 


\hline 


3 &   $10^4$ &   1439  &   1613  &   1191 &    1098 \\ 


\hline 


4 &   $10^4$ &   4935 &    4483 &    2194 &    2138 \\ 


\hline 


1 &   $10^5$ &   11604 &   11896  &  11467 &   10504 \\ 


\hline 


2 &   $10^5$ &   11276 &   11258 &   8879 &    8124 \\ 


\hline 


3 &   $10^5$ &   10958 &   13898 &   8980 &    8165 \\ 


\hline 


4 &   $10^5$ &   42330 &   41963 &   17976 &   17661 \\ 


\hline 


\end{tabular} 


}


\hfill}


\bigskip


\par


Как и предполагалось, на скорость сходимости сильное влияние оказывает 


коэффициент несимметрии системы $k=Pe\,h/2$ $(h$ --- 


шаг по пространству), а не $Pe$ и $h$ в 


отдельности. Установлено, что при различных $Pe$ и $h,$ для которых $k$ 


оставался 


постоянным, значения $\tau_{\text{opt}}$ и $n$ 


(число итераций) практически не изменяются. 


 


Исследовалось влияние коэффициента $k$ и итерационного параметра $\tau$ на 


величину $n$ --- числа итераций, а также на сходимость методов. Численно 


подтвердилось существование таких значений $\tau_{\text{opt}},$ для которых 


число итераций для достижения заданной точности минимально. 


 


Сравнение предложенных методов с методом верхней релаксации и Зейделя 


показало, что при $0<k<1$ последние более эффективны, но при $k>1$ 


предпогалаемые 


методы становятся предпочтительней. Метод Зейделя перестает считать уже при 


$k=1{.}2,$ 


а метод верхней релаксации считает хуже метода нижней релаксации, который 


перестает сходиться при $k>10^3.$ Предлагаемые в данной статье методы не 


только сходятся быстрее вышеупомянутых, но и продолжают работать при 


значениях $k>10^3.$ 


 


Таким образом, рассмотренный выше класс итерационных методов 


предназначен для решения СЛАУ с несамосопряженными положительными матрицами, 


у которых кососимметричная часть является преобладающей. За счет 


использования на верхнем слое оператора $B$ (или, что эквивалентно, 


переобуславливателя), содержащего кососимметричную часть исходного оператора, 


удается решать задачи, в которых число кососимметрии является величиной 


порядка $10^6$ и более. 


 


Автор выражает благодарность А.Л.Чикину за проведенные расчеты. 


 


\begin{thebibliography}{99} 


 


\bibitem{Lan} 


Ландау Л.Д., Лифшиц Е.М. {\em Теоретическая физика.} Т.\,IV. {\em 


Гидродинамика.} -- М.: Наука, 1988 -- 736 с. 





\bibitem{Ro} Роуч П. {\em Вычислительная гидродинамика.} -- М.: Мир, 1980. -- 616 с. 





\bibitem{Neu} 


Neumann M., Plemmons R.J. {\em $M$-matrix characterizations} II: {\em General 


$M$-matrices} 


// Linear and Multilin. Algebra. -- 1980. -- V.\,9. -- \No\,3. -- P.\,211--225. 





\bibitem{Y} 


Young D.M. {\em Iterative solution of large linear systems.} -- New York and 


London: Academic Press, 1971. -- 589 p. 





\bibitem{Mei} 


Meijrlink J.A., Van der Vorst H.A. {\em An iterative solution method for linear 


systems of which the coefficient matrix is a symmetric $M$-matrix} // Math. 


Comput. -- 1977. -- V.\,31. -- \No\,137. -- P.\,148--162.





\bibitem{Vab} 


Вабищевич П.Н. {\em Разностные схемы с центральными разностями для задач 


конвекции-диффузии.} -- М.: Ин-т Мат. Мод. РАН. -- 1993. -- \No\,17. -- 16 с. 





\bibitem{Kru1} 


Крукиер Л.А. {\em Неявные разностные схемы и итерационный метод их решения 


для одного класса систем квазилинейных уравнений} // Изв. вузов. Математика. 


-- 1979. -- \No\,7. -- С.\,41--52. 





\bibitem{Sam} 


Самарский А.А., Николаев Е.С. {\em Методы решения сеточных уравнений.} 


Учеб. пособие. -- М.: Наука, 1978. -- 590 с. 





\bibitem{Hai} 


Хайгеман Л., Янг Д. {\em Прикладные итерационные методы.} -- М.: Мир, 1986. 


-- 448 с. 





\bibitem{Ha} 


Hackbush W. {\em Iterative solution of large sparse systems of equations.} 


-- Berlin.: Springer-Verlag, 1994. -- 429 p. 





\bibitem{Gul} 


Гулин А.В. {\em Устойчивость разностных схем и операторные неравенства} // 


Дифференц. уравнения. -- 1979. -- Т.\,15. -- \No\,12. -- С.\,2238--2250. 





\bibitem{Kru2} 


Крукиер Л.А. {\em О некоторых способах построения оператора $B$ в неявных 


двухслойных итерационных схемах, обеспечивающего их сходимость в случае 


диссипативности оператора $A$} // Изв. вузов. Математика. -- 1983. -- \No\,5. 


-- С.\,41--47. 





\bibitem{Kru3} 


Крукиер Л.А. {\em Математическое моделирование гидродинамики Азовского 


моря при реализации проектов реконструкции его экосистемы} // Матем. моделир. 


-- 1991. -- Т.\,3. -- \No\,9. -- С.\,3--20. 





\bibitem{Kru4} 


Крукиер Л.А., Чикина Л.Г. {\em Решение стационарного уравнения 


конвекции-диффузии в несжимаемых средах с преобладающей конвекцией 


итерационными методами} // Тр. Международн. конф. ``Применение математического 


моделирования для решения задач в науке и технике'', ИММ РАН, ИПМ Уро РАН. -- 


Ижевск, 1996. -- С.\,190--201. 


 


\end{thebibliography} 


 


\vskip 2cm 


 


\post{\it Вычислительный центр}{\it Поступила} 


\post{\it Ростовского государственного университета}{01.11.1996} 


 


 


\end{document} 


